\documentclass[11pt]{article}
%\setlength\textwidth{190mm} \setlength\textheight{237mm}%\setlength\voffset{-18mm} \setlength\hoffset{-33mm}\usepackage[T2A]{fontenc}
\usepackage[utf8]{inputenc}
\usepackage[ukrainian]{babel}
\usepackage{graphicx}
\usepackage[Export]{adjustbox}
\adjustboxset{max size={1.0\linewidth}{1.0\paperheight}}
\graphicspath{{data/pages/}} 
\usepackage{geometry}
\geometry{top=1.5cm,bottom=1.5cm,left=1.2cm,right=1.2cm}
\pagestyle{empty}\setlength\parindent{0mm}
\begin{document}
{\Large \center  Варіанти завдань до програмного проєкту №3\par Прикладна математика, 2023/24 н.р.\par \bigskip\bigskip}

%begin
{\Large \center  Варіант  1\par \bigskip\bigskip}

За квадратною матрицею, заданою об'єктом типу numpy.ndarray, надрукувати послідовність її елементів обходом нижнього чверть-трикутника матриці спіраллю від центрального кута за годинниковою стрілкою.
%z05 S p09 n-

\vspace{2mm}
Якщо попередні вимоги залишають неоднозначність, то перевагу віддавати тому елементу матриці з наявних альтернатив, що має лексикографічно менші індекси.

Вважати, що лівий верхній кут матриці має координати~$(0,0)$.



%end
%begin
{\Large \center  Варіант  2\par \bigskip\bigskip}

За квадратною матрицею, заданою об'єктом типу numpy.ndarray, надрукувати послідовність її елементів обходом нижнього чверть-трикутника матриці ``змійкою'' від нижнього правого кута матриці діагоналями (паралельно головній діагоналі).
%z05 Z p06 n++

\vspace{2mm}
Якщо попередні вимоги залишають неоднозначність, то перевагу віддавати тому елементу матриці з наявних альтернатив, що має лексикографічно менші індекси.

Вважати, що лівий верхній кут матриці має координати~$(0,0)$.



%end
%begin
{\Large \center  Варіант  3\par \bigskip\bigskip}

За квадратною матрицею, заданою об'єктом типу numpy.ndarray, надрукувати послідовність її елементів обходом трикутника над головною діагоналлю матриці ``змійкою'' від верхнього лівого кута матриці горизонталями.
%z01 Z p07 n0+

\vspace{2mm}
Якщо попередні вимоги залишають неоднозначність, то перевагу віддавати тому елементу матриці з наявних альтернатив, що має лексикографічно менші індекси.

Вважати, що лівий верхній кут матриці має координати~$(0,0)$.



%end
%begin
{\Large \center  Варіант  4\par \bigskip\bigskip}

За квадратною матрицею, заданою об'єктом типу numpy.ndarray, надрукувати послідовність її елементів обходом трикутника під головною діагоналлю матриці ``змійкою'' від нижнього лівого кута матриці діагоналями (перпендикулярно головній діагоналі).
%z02 Z p05 n+-

\vspace{2mm}
Якщо попередні вимоги залишають неоднозначність, то перевагу віддавати тому елементу матриці з наявних альтернатив, що має лексикографічно менші індекси.

Вважати, що лівий верхній кут матриці має координати~$(0,0)$.



%end
%begin
{\Large \center  Варіант  5\par \bigskip\bigskip}

За квадратною матрицею, заданою об'єктом типу numpy.ndarray, надрукувати послідовність її елементів обходом центрального ромба матриці спіраллю від його верхнього кута проти годинникової стрілки.
%z00 S p04 n+

\vspace{2mm}
Якщо попередні вимоги залишають неоднозначність, то перевагу віддавати тому елементу матриці з наявних альтернатив, що має лексикографічно менші індекси.

Вважати, що лівий верхній кут матриці має координати~$(0,0)$.



%end
%begin
{\Large \center  Варіант  6\par \bigskip\bigskip}

За квадратною матрицею, заданою об'єктом типу numpy.ndarray, надрукувати послідовність її елементів обходом трикутника під побічною діагоналлю матриці ``змійкою'' від нижнього лівого кута матриці вертикалями.
%z08 Z p05 n+0

\vspace{2mm}
Якщо попередні вимоги залишають неоднозначність, то перевагу віддавати тому елементу матриці з наявних альтернатив, що має лексикографічно менші індекси.

Вважати, що лівий верхній кут матриці має координати~$(0,0)$.



%end
%begin
{\Large \center  Варіант  7\par \bigskip\bigskip}

За квадратною матрицею, заданою об'єктом типу numpy.ndarray, надрукувати послідовність її елементів обходом трикутника під побічною діагоналлю матриці ``змійкою'' від нижнього правого кута матриці вертикалями.
%z08 Z p06 n+0

\vspace{2mm}
Якщо попередні вимоги залишають неоднозначність, то перевагу віддавати тому елементу матриці з наявних альтернатив, що має лексикографічно менші індекси.

Вважати, що лівий верхній кут матриці має координати~$(0,0)$.



%end
%begin
{\Large \center  Варіант  8\par \bigskip\bigskip}

За квадратною матрицею, заданою об'єктом типу numpy.ndarray, надрукувати послідовність її елементів обходом верхнього чверть-трикутника матриці ``змійкою'' від центрального кута горизонталями.
%z03 Z p09 n0+

\vspace{2mm}
Якщо попередні вимоги залишають неоднозначність, то перевагу віддавати тому елементу матриці з наявних альтернатив, що має лексикографічно менші індекси.

Вважати, що лівий верхній кут матриці має координати~$(0,0)$.



%end
%begin
{\Large \center  Варіант  9\par \bigskip\bigskip}

За квадратною матрицею, заданою об'єктом типу numpy.ndarray, надрукувати послідовність її елементів обходом центрального ромба матриці спіраллю від його правого кута за годинниковою стрілкою.
%z00 S p01 n-

\vspace{2mm}
Якщо попередні вимоги залишають неоднозначність, то перевагу віддавати тому елементу матриці з наявних альтернатив, що має лексикографічно менші індекси.

Вважати, що лівий верхній кут матриці має координати~$(0,0)$.



%end
%begin
{\Large \center  Варіант  10\par \bigskip\bigskip}

За квадратною матрицею, заданою об'єктом типу numpy.ndarray, надрукувати послідовність її елементів обходом трикутника над побічною діагоналлю матриці ``змійкою'' від верхнього правого кута матриці діагоналями (паралельно головній діагоналі).
%z07 Z p08 n++

\vspace{2mm}
Якщо попередні вимоги залишають неоднозначність, то перевагу віддавати тому елементу матриці з наявних альтернатив, що має лексикографічно менші індекси.

Вважати, що лівий верхній кут матриці має координати~$(0,0)$.



%end
\newpage
%begin
{\Large \center  Варіант  11\par \bigskip\bigskip}

За квадратною матрицею, заданою об'єктом типу numpy.ndarray, надрукувати послідовність її елементів обходом лівого чверть-трикутника матриці спіраллю від нижнього лівого кута матриці проти годинникової стрілки.
%z06 S p05 n+

\vspace{2mm}
Якщо попередні вимоги залишають неоднозначність, то перевагу віддавати тому елементу матриці з наявних альтернатив, що має лексикографічно менші індекси.

Вважати, що лівий верхній кут матриці має координати~$(0,0)$.



%end
%begin
{\Large \center  Варіант  12\par \bigskip\bigskip}

За квадратною матрицею, заданою об'єктом типу numpy.ndarray, надрукувати послідовність її елементів обходом трикутника над побічною діагоналлю матриці ``змійкою'' від нижнього лівого кута матриці діагоналями (перпендикулярно головній діагоналі).
%z07 Z p05 n+-

\vspace{2mm}
Якщо попередні вимоги залишають неоднозначність, то перевагу віддавати тому елементу матриці з наявних альтернатив, що має лексикографічно менші індекси.

Вважати, що лівий верхній кут матриці має координати~$(0,0)$.



%end
%begin
{\Large \center  Варіант  13\par \bigskip\bigskip}

За квадратною матрицею, заданою об'єктом типу numpy.ndarray, надрукувати послідовність її елементів обходом верхнього чверть-трикутника матриці ``змійкою'' від верхнього правого кута матриці горизонталями.
%z03 Z p08 n0+

\vspace{2mm}
Якщо попередні вимоги залишають неоднозначність, то перевагу віддавати тому елементу матриці з наявних альтернатив, що має лексикографічно менші індекси.

Вважати, що лівий верхній кут матриці має координати~$(0,0)$.



%end
%begin
{\Large \center  Варіант  14\par \bigskip\bigskip}

За квадратною матрицею, заданою об'єктом типу numpy.ndarray, надрукувати послідовність її елементів обходом верхнього чверть-трикутника матриці ``змійкою'' від верхнього лівого кута матриці діагоналями (паралельно головній діагоналі).
%z03 Z p07 n++

\vspace{2mm}
Якщо попередні вимоги залишають неоднозначність, то перевагу віддавати тому елементу матриці з наявних альтернатив, що має лексикографічно менші індекси.

Вважати, що лівий верхній кут матриці має координати~$(0,0)$.



%end
%begin
{\Large \center  Варіант  15\par \bigskip\bigskip}

За квадратною матрицею, заданою об'єктом типу numpy.ndarray, надрукувати послідовність її елементів обходом центрального ромба матриці ``змійкою'' від його нижнього кута діагоналями (перпендикулярно головній діагоналі).
%z00 Z p03 n+-

\vspace{2mm}
Якщо попередні вимоги залишають неоднозначність, то перевагу віддавати тому елементу матриці з наявних альтернатив, що має лексикографічно менші індекси.

Вважати, що лівий верхній кут матриці має координати~$(0,0)$.



%end
%begin
{\Large \center  Варіант  16\par \bigskip\bigskip}

За квадратною матрицею, заданою об'єктом типу numpy.ndarray, надрукувати послідовність її елементів обходом центрального ромба матриці спіраллю від його лівого кута за годинниковою стрілкою.
%z00 S p02 n-

\vspace{2mm}
Якщо попередні вимоги залишають неоднозначність, то перевагу віддавати тому елементу матриці з наявних альтернатив, що має лексикографічно менші індекси.

Вважати, що лівий верхній кут матриці має координати~$(0,0)$.



%end
%begin
{\Large \center  Варіант  17\par \bigskip\bigskip}

За квадратною матрицею, заданою об'єктом типу numpy.ndarray, надрукувати послідовність її елементів обходом лівого чверть-трикутника матриці ``змійкою'' від центрального кута вертикалями.
%z06 Z p09 n+0

\vspace{2mm}
Якщо попередні вимоги залишають неоднозначність, то перевагу віддавати тому елементу матриці з наявних альтернатив, що має лексикографічно менші індекси.

Вважати, що лівий верхній кут матриці має координати~$(0,0)$.



%end
%begin
{\Large \center  Варіант  18\par \bigskip\bigskip}

За квадратною матрицею, заданою об'єктом типу numpy.ndarray, надрукувати послідовність її елементів обходом трикутника під побічною діагоналлю матриці спіраллю від верхнього правого кута матриці проти годинникової стрілки.
%z08 S p08 n+

\vspace{2mm}
Якщо попередні вимоги залишають неоднозначність, то перевагу віддавати тому елементу матриці з наявних альтернатив, що має лексикографічно менші індекси.

Вважати, що лівий верхній кут матриці має координати~$(0,0)$.



%end
%begin
{\Large \center  Варіант  19\par \bigskip\bigskip}

За квадратною матрицею, заданою об'єктом типу numpy.ndarray, надрукувати послідовність її елементів обходом трикутника над головною діагоналлю матриці ``змійкою'' від нижнього правого кута матриці діагоналями (перпендикулярно головній діагоналі).
%z01 Z p06 n+-

\vspace{2mm}
Якщо попередні вимоги залишають неоднозначність, то перевагу віддавати тому елементу матриці з наявних альтернатив, що має лексикографічно менші індекси.

Вважати, що лівий верхній кут матриці має координати~$(0,0)$.



%end
%begin
{\Large \center  Варіант  20\par \bigskip\bigskip}

За квадратною матрицею, заданою об'єктом типу numpy.ndarray, надрукувати послідовність її елементів обходом правого чверть-трикутника матриці спіраллю від верхнього правого кута матриці за годинниковою стрілкою.
%z04 S p08 n-

\vspace{2mm}
Якщо попередні вимоги залишають неоднозначність, то перевагу віддавати тому елементу матриці з наявних альтернатив, що має лексикографічно менші індекси.

Вважати, що лівий верхній кут матриці має координати~$(0,0)$.



%end
%begin
{\Large \center  Варіант  21\par \bigskip\bigskip}

За квадратною матрицею, заданою об'єктом типу numpy.ndarray, надрукувати послідовність її елементів обходом трикутника над головною діагоналлю матриці ``змійкою'' від верхнього правого кута матриці діагоналями (паралельно головній діагоналі).
%z01 Z p08 n++

\vspace{2mm}
Якщо попередні вимоги залишають неоднозначність, то перевагу віддавати тому елементу матриці з наявних альтернатив, що має лексикографічно менші індекси.

Вважати, що лівий верхній кут матриці має координати~$(0,0)$.



%end
%begin
{\Large \center  Варіант  22\par \bigskip\bigskip}

За квадратною матрицею, заданою об'єктом типу numpy.ndarray, надрукувати послідовність її елементів обходом правого чверть-трикутника матриці ``змійкою'' від нижнього правого кута матриці вертикалями.
%z04 Z p06 n+0

\vspace{2mm}
Якщо попередні вимоги залишають неоднозначність, то перевагу віддавати тому елементу матриці з наявних альтернатив, що має лексикографічно менші індекси.

Вважати, що лівий верхній кут матриці має координати~$(0,0)$.



%end
%begin
{\Large \center  Варіант  23\par \bigskip\bigskip}

За квадратною матрицею, заданою об'єктом типу numpy.ndarray, надрукувати послідовність її елементів обходом трикутника над побічною діагоналлю матриці ``змійкою'' від верхнього лівого кута матриці горизонталями.
%z07 Z p07 n0+

\vspace{2mm}
Якщо попередні вимоги залишають неоднозначність, то перевагу віддавати тому елементу матриці з наявних альтернатив, що має лексикографічно менші індекси.

Вважати, що лівий верхній кут матриці має координати~$(0,0)$.



%end
%begin
{\Large \center  Варіант  24\par \bigskip\bigskip}

За квадратною матрицею, заданою об'єктом типу numpy.ndarray, надрукувати послідовність її елементів обходом правого чверть-трикутника матриці спіраллю від центрального кута проти годинникової стрілки.
%z04 S p09 n+

\vspace{2mm}
Якщо попередні вимоги залишають неоднозначність, то перевагу віддавати тому елементу матриці з наявних альтернатив, що має лексикографічно менші індекси.

Вважати, що лівий верхній кут матриці має координати~$(0,0)$.



%end
%begin
{\Large \center  Варіант  25\par \bigskip\bigskip}

За квадратною матрицею, заданою об'єктом типу numpy.ndarray, надрукувати послідовність її елементів обходом трикутника під головною діагоналлю матриці ``змійкою'' від нижнього правого кута матриці діагоналями (перпендикулярно головній діагоналі).
%z02 Z p06 n+-

\vspace{2mm}
Якщо попередні вимоги залишають неоднозначність, то перевагу віддавати тому елементу матриці з наявних альтернатив, що має лексикографічно менші індекси.

Вважати, що лівий верхній кут матриці має координати~$(0,0)$.



%end
%begin
{\Large \center  Варіант  26\par \bigskip\bigskip}

За квадратною матрицею, заданою об'єктом типу numpy.ndarray, надрукувати послідовність її елементів обходом трикутника під головною діагоналлю матриці ``змійкою'' від верхнього лівого кута матриці горизонталями.
%z02 Z p07 n0+

\vspace{2mm}
Якщо попередні вимоги залишають неоднозначність, то перевагу віддавати тому елементу матриці з наявних альтернатив, що має лексикографічно менші індекси.

Вважати, що лівий верхній кут матриці має координати~$(0,0)$.



%end
%begin
{\Large \center  Варіант  27\par \bigskip\bigskip}

За квадратною матрицею, заданою об'єктом типу numpy.ndarray, надрукувати послідовність її елементів обходом нижнього чверть-трикутника матриці ``змійкою'' від нижнього лівого кута матриці вертикалями.
%z05 Z p05 n+0

\vspace{2mm}
Якщо попередні вимоги залишають неоднозначність, то перевагу віддавати тому елементу матриці з наявних альтернатив, що має лексикографічно менші індекси.

Вважати, що лівий верхній кут матриці має координати~$(0,0)$.



%end
%begin
{\Large \center  Варіант  28\par \bigskip\bigskip}

За квадратною матрицею, заданою об'єктом типу numpy.ndarray, надрукувати послідовність її елементів обходом лівого чверть-трикутника матриці спіраллю від верхнього лівого кута матриці проти годинникової стрілки.
%z06 S p07 n+

\vspace{2mm}
Якщо попередні вимоги залишають неоднозначність, то перевагу віддавати тому елементу матриці з наявних альтернатив, що має лексикографічно менші індекси.

Вважати, що лівий верхній кут матриці має координати~$(0,0)$.



%end
%begin
{\Large \center  Варіант  29\par \bigskip\bigskip}

За квадратною матрицею, заданою об'єктом типу numpy.ndarray, надрукувати послідовність її елементів обходом трикутника під побічною діагоналлю матриці ``змійкою'' від верхнього правого кута матриці діагоналями (паралельно головній діагоналі).
%z08 Z p08 n++

\vspace{2mm}
Якщо попередні вимоги залишають неоднозначність, то перевагу віддавати тому елементу матриці з наявних альтернатив, що має лексикографічно менші індекси.

Вважати, що лівий верхній кут матриці має координати~$(0,0)$.



%end
%begin
{\Large \center  Варіант  30\par \bigskip\bigskip}

За квадратною матрицею, заданою об'єктом типу numpy.ndarray, надрукувати послідовність її елементів обходом лівого чверть-трикутника матриці спіраллю від центрального кута за годинниковою стрілкою.
%z06 S p09 n-

\vspace{2mm}
Якщо попередні вимоги залишають неоднозначність, то перевагу віддавати тому елементу матриці з наявних альтернатив, що має лексикографічно менші індекси.

Вважати, що лівий верхній кут матриці має координати~$(0,0)$.



%end
%begin
{\Large \center  Варіант  31\par \bigskip\bigskip}

За квадратною матрицею, заданою об'єктом типу numpy.ndarray, надрукувати послідовність її елементів обходом трикутника над головною діагоналлю матриці спіраллю від нижнього правого кута матриці проти годинникової стрілки.
%z01 S p06 n+

\vspace{2mm}
Якщо попередні вимоги залишають неоднозначність, то перевагу віддавати тому елементу матриці з наявних альтернатив, що має лексикографічно менші індекси.

Вважати, що лівий верхній кут матриці має координати~$(0,0)$.



%end
%begin
{\Large \center  Варіант  32\par \bigskip\bigskip}

За квадратною матрицею, заданою об'єктом типу numpy.ndarray, надрукувати послідовність її елементів обходом трикутника над головною діагоналлю матриці ``змійкою'' від верхнього лівого кута матриці горизонталями.
%z01 Z p07 n0+

\vspace{2mm}
Якщо попередні вимоги залишають неоднозначність, то перевагу віддавати тому елементу матриці з наявних альтернатив, що має лексикографічно менші індекси.

Вважати, що лівий верхній кут матриці має координати~$(0,0)$.



%end
%begin
{\Large \center  Варіант  33\par \bigskip\bigskip}

За квадратною матрицею, заданою об'єктом типу numpy.ndarray, надрукувати послідовність її елементів обходом правого чверть-трикутника матриці ``змійкою'' від верхнього правого кута матриці діагоналями (перпендикулярно головній діагоналі).
%z04 Z p08 n+-

\vspace{2mm}
Якщо попередні вимоги залишають неоднозначність, то перевагу віддавати тому елементу матриці з наявних альтернатив, що має лексикографічно менші індекси.

Вважати, що лівий верхній кут матриці має координати~$(0,0)$.



%end
%begin
{\Large \center  Варіант  34\par \bigskip\bigskip}

За квадратною матрицею, заданою об'єктом типу numpy.ndarray, надрукувати послідовність її елементів обходом нижнього чверть-трикутника матриці ``змійкою'' від нижнього лівого кута матриці вертикалями.
%z05 Z p05 n+0

\vspace{2mm}
Якщо попередні вимоги залишають неоднозначність, то перевагу віддавати тому елементу матриці з наявних альтернатив, що має лексикографічно менші індекси.

Вважати, що лівий верхній кут матриці має координати~$(0,0)$.



%end
%begin
{\Large \center  Варіант  35\par \bigskip\bigskip}

За квадратною матрицею, заданою об'єктом типу numpy.ndarray, надрукувати послідовність її елементів обходом нижнього чверть-трикутника матриці спіраллю від центрального кута за годинниковою стрілкою.
%z05 S p09 n-

\vspace{2mm}
Якщо попередні вимоги залишають неоднозначність, то перевагу віддавати тому елементу матриці з наявних альтернатив, що має лексикографічно менші індекси.

Вважати, що лівий верхній кут матриці має координати~$(0,0)$.



%end
%begin
{\Large \center  Варіант  36\par \bigskip\bigskip}

За квадратною матрицею, заданою об'єктом типу numpy.ndarray, надрукувати послідовність її елементів обходом трикутника під побічною діагоналлю матриці ``змійкою'' від нижнього правого кута матриці діагоналями (паралельно головній діагоналі).
%z08 Z p06 n++

\vspace{2mm}
Якщо попередні вимоги залишають неоднозначність, то перевагу віддавати тому елементу матриці з наявних альтернатив, що має лексикографічно менші індекси.

Вважати, що лівий верхній кут матриці має координати~$(0,0)$.



%end
%begin
{\Large \center  Варіант  37\par \bigskip\bigskip}

За квадратною матрицею, заданою об'єктом типу numpy.ndarray, надрукувати послідовність її елементів обходом верхнього чверть-трикутника матриці спіраллю від центрального кута за годинниковою стрілкою.
%z03 S p09 n-

\vspace{2mm}
Якщо попередні вимоги залишають неоднозначність, то перевагу віддавати тому елементу матриці з наявних альтернатив, що має лексикографічно менші індекси.

Вважати, що лівий верхній кут матриці має координати~$(0,0)$.



%end
%begin
{\Large \center  Варіант  38\par \bigskip\bigskip}

За квадратною матрицею, заданою об'єктом типу numpy.ndarray, надрукувати послідовність її елементів обходом нижнього чверть-трикутника матриці ``змійкою'' від нижнього правого кута матриці вертикалями.
%z05 Z p06 n+0

\vspace{2mm}
Якщо попередні вимоги залишають неоднозначність, то перевагу віддавати тому елементу матриці з наявних альтернатив, що має лексикографічно менші індекси.

Вважати, що лівий верхній кут матриці має координати~$(0,0)$.



%end
%begin
{\Large \center  Варіант  39\par \bigskip\bigskip}

За квадратною матрицею, заданою об'єктом типу numpy.ndarray, надрукувати послідовність її елементів обходом трикутника під головною діагоналлю матриці ``змійкою'' від верхнього лівого кута матриці діагоналями (паралельно головній діагоналі).
%z02 Z p07 n++

\vspace{2mm}
Якщо попередні вимоги залишають неоднозначність, то перевагу віддавати тому елементу матриці з наявних альтернатив, що має лексикографічно менші індекси.

Вважати, що лівий верхній кут матриці має координати~$(0,0)$.



%end
%begin
{\Large \center  Варіант  40\par \bigskip\bigskip}

За квадратною матрицею, заданою об'єктом типу numpy.ndarray, надрукувати послідовність її елементів обходом правого чверть-трикутника матриці ``змійкою'' від нижнього правого кута матриці діагоналями (перпендикулярно головній діагоналі).
%z04 Z p06 n+-

\vspace{2mm}
Якщо попередні вимоги залишають неоднозначність, то перевагу віддавати тому елементу матриці з наявних альтернатив, що має лексикографічно менші індекси.

Вважати, що лівий верхній кут матриці має координати~$(0,0)$.



%end
%begin
{\Large \center  Варіант  41\par \bigskip\bigskip}

За квадратною матрицею, заданою об'єктом типу numpy.ndarray, надрукувати послідовність її елементів обходом трикутника над побічною діагоналлю матриці спіраллю від верхнього правого кута матриці проти годинникової стрілки.
%z07 S p08 n+

\vspace{2mm}
Якщо попередні вимоги залишають неоднозначність, то перевагу віддавати тому елементу матриці з наявних альтернатив, що має лексикографічно менші індекси.

Вважати, що лівий верхній кут матриці має координати~$(0,0)$.



%end
%begin
{\Large \center  Варіант  42\par \bigskip\bigskip}

За квадратною матрицею, заданою об'єктом типу numpy.ndarray, надрукувати послідовність її елементів обходом верхнього чверть-трикутника матриці ``змійкою'' від верхнього лівого кута матриці горизонталями.
%z03 Z p07 n0+

\vspace{2mm}
Якщо попередні вимоги залишають неоднозначність, то перевагу віддавати тому елементу матриці з наявних альтернатив, що має лексикографічно менші індекси.

Вважати, що лівий верхній кут матриці має координати~$(0,0)$.



%end
%begin
{\Large \center  Варіант  43\par \bigskip\bigskip}

За квадратною матрицею, заданою об'єктом типу numpy.ndarray, надрукувати послідовність її елементів обходом трикутника над побічною діагоналлю матриці ``змійкою'' від нижнього лівого кута матриці горизонталями.
%z07 Z p05 n0+

\vspace{2mm}
Якщо попередні вимоги залишають неоднозначність, то перевагу віддавати тому елементу матриці з наявних альтернатив, що має лексикографічно менші індекси.

Вважати, що лівий верхній кут матриці має координати~$(0,0)$.



%end
%begin
{\Large \center  Варіант  44\par \bigskip\bigskip}

За квадратною матрицею, заданою об'єктом типу numpy.ndarray, надрукувати послідовність її елементів обходом центрального ромба матриці спіраллю від його верхнього кута проти годинникової стрілки.
%z00 S p04 n+

\vspace{2mm}
Якщо попередні вимоги залишають неоднозначність, то перевагу віддавати тому елементу матриці з наявних альтернатив, що має лексикографічно менші індекси.

Вважати, що лівий верхній кут матриці має координати~$(0,0)$.



%end
%begin
{\Large \center  Варіант  45\par \bigskip\bigskip}

За квадратною матрицею, заданою об'єктом типу numpy.ndarray, надрукувати послідовність її елементів обходом трикутника під головною діагоналлю матриці спіраллю від нижнього лівого кута матриці за годинниковою стрілкою.
%z02 S p05 n-

\vspace{2mm}
Якщо попередні вимоги залишають неоднозначність, то перевагу віддавати тому елементу матриці з наявних альтернатив, що має лексикографічно менші індекси.

Вважати, що лівий верхній кут матриці має координати~$(0,0)$.



%end
%begin
{\Large \center  Варіант  46\par \bigskip\bigskip}

За квадратною матрицею, заданою об'єктом типу numpy.ndarray, надрукувати послідовність її елементів обходом правого чверть-трикутника матриці ``змійкою'' від центрального кута діагоналями (перпендикулярно головній діагоналі).
%z04 Z p09 n+-

\vspace{2mm}
Якщо попередні вимоги залишають неоднозначність, то перевагу віддавати тому елементу матриці з наявних альтернатив, що має лексикографічно менші індекси.

Вважати, що лівий верхній кут матриці має координати~$(0,0)$.



%end
%begin
{\Large \center  Варіант  47\par \bigskip\bigskip}

За квадратною матрицею, заданою об'єктом типу numpy.ndarray, надрукувати послідовність її елементів обходом центрального ромба матриці ``змійкою'' від його нижнього кута вертикалями.
%z00 Z p03 n+0

\vspace{2mm}
Якщо попередні вимоги залишають неоднозначність, то перевагу віддавати тому елементу матриці з наявних альтернатив, що має лексикографічно менші індекси.

Вважати, що лівий верхній кут матриці має координати~$(0,0)$.



%end
%begin
{\Large \center  Варіант  48\par \bigskip\bigskip}

За квадратною матрицею, заданою об'єктом типу numpy.ndarray, надрукувати послідовність її елементів обходом трикутника під головною діагоналлю матриці ``змійкою'' від нижнього правого кута матриці діагоналями (паралельно головній діагоналі).
%z02 Z p06 n++

\vspace{2mm}
Якщо попередні вимоги залишають неоднозначність, то перевагу віддавати тому елементу матриці з наявних альтернатив, що має лексикографічно менші індекси.

Вважати, що лівий верхній кут матриці має координати~$(0,0)$.



%end
%begin
{\Large \center  Варіант  49\par \bigskip\bigskip}

За квадратною матрицею, заданою об'єктом типу numpy.ndarray, надрукувати послідовність її елементів обходом верхнього чверть-трикутника матриці ``змійкою'' від верхнього правого кута матриці горизонталями.
%z03 Z p08 n0+

\vspace{2mm}
Якщо попередні вимоги залишають неоднозначність, то перевагу віддавати тому елементу матриці з наявних альтернатив, що має лексикографічно менші індекси.

Вважати, що лівий верхній кут матриці має координати~$(0,0)$.



%end
%begin
{\Large \center  Варіант  50\par \bigskip\bigskip}

За квадратною матрицею, заданою об'єктом типу numpy.ndarray, надрукувати послідовність її елементів обходом центрального ромба матриці спіраллю від його лівого кута проти годинникової стрілки.
%z00 S p02 n+

\vspace{2mm}
Якщо попередні вимоги залишають неоднозначність, то перевагу віддавати тому елементу матриці з наявних альтернатив, що має лексикографічно менші індекси.

Вважати, що лівий верхній кут матриці має координати~$(0,0)$.



%end
%begin
{\Large \center  Варіант  51\par \bigskip\bigskip}

За квадратною матрицею, заданою об'єктом типу numpy.ndarray, надрукувати послідовність її елементів обходом трикутника під побічною діагоналлю матриці спіраллю від нижнього лівого кута матриці за годинниковою стрілкою.
%z08 S p05 n-

\vspace{2mm}
Якщо попередні вимоги залишають неоднозначність, то перевагу віддавати тому елементу матриці з наявних альтернатив, що має лексикографічно менші індекси.

Вважати, що лівий верхній кут матриці має координати~$(0,0)$.



%end
%begin
{\Large \center  Варіант  52\par \bigskip\bigskip}

За квадратною матрицею, заданою об'єктом типу numpy.ndarray, надрукувати послідовність її елементів обходом трикутника над головною діагоналлю матриці ``змійкою'' від верхнього правого кута матриці діагоналями (перпендикулярно головній діагоналі).
%z01 Z p08 n+-

\vspace{2mm}
Якщо попередні вимоги залишають неоднозначність, то перевагу віддавати тому елементу матриці з наявних альтернатив, що має лексикографічно менші індекси.

Вважати, що лівий верхній кут матриці має координати~$(0,0)$.



%end
%begin
{\Large \center  Варіант  53\par \bigskip\bigskip}

За квадратною матрицею, заданою об'єктом типу numpy.ndarray, надрукувати послідовність її елементів обходом лівого чверть-трикутника матриці ``змійкою'' від верхнього лівого кута матриці вертикалями.
%z06 Z p07 n+0

\vspace{2mm}
Якщо попередні вимоги залишають неоднозначність, то перевагу віддавати тому елементу матриці з наявних альтернатив, що має лексикографічно менші індекси.

Вважати, що лівий верхній кут матриці має координати~$(0,0)$.



%end
%begin
{\Large \center  Варіант  54\par \bigskip\bigskip}

За квадратною матрицею, заданою об'єктом типу numpy.ndarray, надрукувати послідовність її елементів обходом трикутника над побічною діагоналлю матриці ``змійкою'' від верхнього лівого кута матриці діагоналями (паралельно головній діагоналі).
%z07 Z p07 n++

\vspace{2mm}
Якщо попередні вимоги залишають неоднозначність, то перевагу віддавати тому елементу матриці з наявних альтернатив, що має лексикографічно менші індекси.

Вважати, що лівий верхній кут матриці має координати~$(0,0)$.



%end
%begin
{\Large \center  Варіант  55\par \bigskip\bigskip}

За квадратною матрицею, заданою об'єктом типу numpy.ndarray, надрукувати послідовність її елементів обходом лівого чверть-трикутника матриці спіраллю від нижнього лівого кута матриці проти годинникової стрілки.
%z06 S p05 n+

\vspace{2mm}
Якщо попередні вимоги залишають неоднозначність, то перевагу віддавати тому елементу матриці з наявних альтернатив, що має лексикографічно менші індекси.

Вважати, що лівий верхній кут матриці має координати~$(0,0)$.



%end
\end{document}
